% s-209-p209-m003; original PDF 209, 210
\Aphorism{s-209-p209-m003}{302}
Neben dem, was das Kunstwerk für Andere ist, ist es auch noch etwas für seinen Schöpfer: ein \Sperr{Versteck} für lauter Persönlichstes, für Anspielungen und Beziehungen, die kein Anderer versteht, ein Schlupfwinkel für kleine menschliche Besonderheiten und Zufälle, die eigentlich unter der Grenze des Darstellbaren liegen und doch unentbehrlich waren zur Entstehung des „Willens zum Kunstwerk“. Der Künstler allein weiss, wieviel wirkliche Erinnerungen an Land und Volk, an ein schönes Weib oder ein Glas guten Weines er in sein Werk hineingepackt hat, unfindbar für Leser und Zuschauer, rein als Herzenstrost und zärtliches Geheimniss für sich selbst; und wenn er den Besucher vor die Staffelei führt, so lächelt ihn ein Portrait an, wo der Kunstfreund nur eine Landschaft, eine Historie, kurzum reine subjectlose Kunst gewahrt. Dem Künstler flüstert sein Werk etwas ganz Vertrauliches zu, einen Namen, eine Kleinigkeit, irgend etwas Dummes und Erotisches; von dieser Privatbeziehung erfährt der Betrachter nichts — er hat sogar ein ästhetisches Gebot daraus gemacht, nichts davon zu erfahren.
\par

